\section{来源审计与课程主线}

这讲课横跨临床语言、蛋白质序列和基因调控，表面上像四篇论文的串讲，真正的主线却很统一：\textbf{先把生物医学对象写成序列，再让 Transformer 的表示、注意力与生成能力适配每个领域的约束。} 因此，本讲不能只摘录 Med-PaLM 的结果，也不能把后半段压缩成“Transformer 还能用于生物学”。我们需要追问每个案例的输入是什么、监督信号是什么、长上下文或噪声在哪里、模型输出怎样被验证，以及实验结果不能推出什么。

\subsection{本讲如何重建}

课程官网给出录像和三项推荐工作，但没有公开独立 slide deck。本讲以 Stanford Online 的全高清官方录像为视觉母版：画面尺寸为 1920 像素乘 1080 像素；每两秒采样得到 2,045 帧，稳定性与差异过滤留下 103 个高召回候选，人工逐张比较后保留 84 个独立教学状态。所有图像都从全宽原视频重新提取，没有裁掉右侧内容；官方 \texttt{en-US} 人工字幕则用于恢复讲者的动机、限制、课堂问答和过渡。

\begin{longtable}{@{}L{0.18\textwidth}L{0.28\textwidth}L{0.42\textwidth}@{}}
\toprule
证据层 & 本地材料 & 在讲义中的职责 \\
\midrule
官方录像 & 1:08:09 视频与 84 张教学画面 & 决定课堂顺序、表格、公式、结果图和讲者实际展示的证据 \\
官方人工字幕 & \texttt{lecture17.en.srt} 与 timed transcript & 恢复“为什么”、图表纠错、风险解释、工程影响和未来判断 \\
指定论文 & Med-PaLM、ProtNLM、Enformer & 核对任务定义、模型机制、实验对象与结论边界 \\
背景论文 & Performer、DeepConsensus、PaLM、Flan-PaLM & 补齐课堂压缩展示的公式与系统上下文 \\
审计文件 & manifest、coverage、ledger、blueprint & 证明每个视觉节点与 teacher voice 节点都有明确去向 \\
\bottomrule
\end{longtable}

\begin{warningbox}{历史边界：2023 年 2 月 21 日}
本讲重建课堂当日的知识状态。后来的 Med-PaLM 2、Gemini 系列、生物医学多模态模型、监管事件与更新 benchmark 都不倒灌为课堂证据。讲者最后提出的是 foundation biomedical AI 的研究愿景，不是一个已经统一临床、蛋白质和基因组任务的成品系统。
\end{warningbox}

\subsection{三段旅程：临床、蛋白质、基因组}

标题页只给出主题，agenda 页却揭示了课程设计：先从第一性原理回答“为什么是 Transformer”，再依次进入 clinical、proteins、genomics，最后讨论未来。读这两页时要注意，领域顺序大致沿着抽象层次下沉：从医生和患者使用的自然语言，走向氨基酸序列，再走向核苷酸与调控信号；但每向下一层，模型输出都必须重新定义，不能沿用聊天机器人的评价标准。

\lecturefigure{slide-01-title-transformers-biomedicine.jpg}{Biomedical Transformers：课程标题与讲者}{官方录像恢复幻灯片 V001；Stanford CS25 Lecture 17}

\lecturefigure{slide-02-agenda.jpg}{课程路线：临床应用、蛋白质、基因组与未来}{官方录像恢复幻灯片 V002；Stanford CS25 Lecture 17}

\teachervoice{讲者开场说自己的目标是“peel back the curtains”，让学生看到 AI 与 biomedicine 交叉处正在发生的创新。他并不打算证明某一个模型统治所有任务，而是选择跨越 biomedical stack 的案例，让共同机制与领域差异同时显现。}

\begin{importantbox}{贯穿全课的五个问题}
对每个案例都应依次回答：输入序列如何构造？上下文为什么难？监督信号来自哪里？输出如何被人或实验验证？结果支持的是 benchmark 改进、科学发现、临床辅助，还是实际部署？如果这五层混在一起，模型分数很容易被误写成医学能力。
\end{importantbox}

\subsection{本章小结}

本讲的证据链由官方录像、人工字幕与五组主要论文构成。课堂结构不是“医疗聊天机器人合集”，而是比较 Transformer 如何处理不同生物医学序列、如何改变复杂度、如何定义输出，以及如何把模型结果接回临床或实验验证。

